\section*{第 \ref{ch:g10:proba} 章 --- 概率与抽样}

\begin{solution}{exo:g10:proba:1}
共 $6$ 种结果，每种结果的概率都是 $\frac16$。
“掷出 $6$”：$\frac16$。
“奇数”$= \{1, 3, 5\}$：$\frac36 = \frac12$。
“不超过 $4$”$= \{1,2,3,4\}$：$\frac46 = \frac23$。
“掷出 $7$”：不可能事件，概率为 $0$。
\end{solution}

\begin{solution}{exo:g10:proba:2}
共有 $10$ 颗等可能的弹珠。
$\P(\text{绿色}) = \frac{5}{10} = \frac12$。
$\P(\text{非黑色}) = 1 - \frac{2}{10} = \frac{8}{10} = \frac45$。
绿色与黄色互斥：
$\P(\text{绿色或黄色}) = \frac{5}{10} + \frac{3}{10} = \frac{8}{10}
= \frac45$。
\end{solution}

\begin{solution}{exo:g10:proba:3}
设 $S$：“学习西班牙语”，$G$：“学习德语”。则
$\P(S) = \frac{18}{30}$，$\P(G) = \frac{10}{30}$，
$\P(S \cap G) = \frac{4}{30}$。加法公式：
\[
\P(S \cup G) = \frac{18}{30} + \frac{10}{30} - \frac{4}{30}
= \frac{24}{30} = \frac45 .
\]
对立事件：$\P(\text{两门都不学}) = 1 - \frac45 = \frac15$。
\end{solution}

\begin{solution}{exo:g10:proba:4}
树状图中共有四条路径 HH、HT、TH、TT，每条路径的概率均为
$\frac12 \times \frac12 = \frac14$。
两次均为正面：$\frac14$。恰好一次为正面（HT 或 TH）：$\frac24 = \frac12$。
至少一次为正面：$1 - \P(\text{TT}) = 1 - \frac14 = \frac34$。
\end{solution}

\begin{solution}{exo:g10:proba:5}
第一次抽取：抽到红球的概率为 $\frac{4}{10}$，抽到白球的概率为 $\frac{6}{10}$。
第二次抽取\emph{不放回}：若第一次取到红球，则剩下的球中
有 $3$ 个红球、$6$ 个白球，共 $9$ 个；若第一次取到白球，则剩下的球中有 $4$ 个红球、$5$ 个白球，共 $9$ 个。

两次均为红球：$\frac{4}{10} \times \frac39 = \frac{12}{90} = \frac{2}{15}$。

颜色不同：路径 RW 与 WR：
\[
\frac{4}{10} \times \frac69 + \frac{6}{10} \times \frac49
= \frac{24}{90} + \frac{24}{90} = \frac{48}{90} = \frac{8}{15}.
\]
\end{solution}

\begin{solution}{exo:g10:proba:6}
\emph{1.} 每个骰子独立地，每一面出现的概率均为
$\frac16$，因此 $36$ 个有序数对等可能，点数之和
由数对唯一确定。

\emph{2.} 点数之和为 $7$：数对 $(1,6), (2,5), (3,4), (4,3), (5,2), (6,1)$，
故 $\P = \frac{6}{36} = \frac16$。点数之和为 $12$：只有 $(6,6)$，
$\P = \frac{1}{36}$。点数之和不超过 $4$：和为 $2$、$3$ 或 $4$ 的数对：
$(1,1)$；$(1,2), (2,1)$；$(1,3), (2,2), (3,1)$——共六个数对，故
$\P = \frac{6}{36} = \frac16$。

\emph{3.} 各点数之和对应的数对个数先从 $1$（和为 $2$）
增加到 $6$（和为 $7$），然后减少：最可能出现的点数之和是 $7$。
\end{solution}

\begin{solution}{exo:g10:proba:7}
每道题答对的概率均为 $\frac14$，
且相互独立。两级树状图如下：

两题都答对：$\frac14 \times \frac14 = \frac{1}{16}$。

恰好答对一题：路径“对—错”与“错—对”：
$\frac14 \times \frac34 + \frac34 \times \frac14 = \frac{6}{16}
= \frac38$。

两题都答错：$\frac34 \times \frac34 = \frac{9}{16}$。

检验：$\frac{1}{16} + \frac{6}{16} + \frac{9}{16} = 1$。
\end{solution}

\begin{solution}{exo:g10:proba:8}
\emph{1.} 观测频率为 $\frac{396}{900} = 0.44$。取 $n = 900$：
$\frac{1}{\sqrt{900}} = \frac{1}{30} \approx 0.033$，因此波动
区间约为 $\intcc{0.467}{0.533}$。

\emph{2.} 观测值 $0.44$ 落在区间之外：抽取 $900$ 名选民组成样本，
若总体支持率为 $50\%$，样本很少会偏离得这么远。这次
民意调查对所声称的 $50\%$ 提出了严重质疑。
\end{solution}

\begin{solution}{exo:g10:proba:9}
\emph{1.} $\P(\text{赢得 } 6) = \frac16$（掷出 $6$），
$\P(\text{赢得 } 3) = \frac16$（掷出 $5$），
$\P(\text{赢得 } 0) = \frac46 = \frac23$。

\emph{2.} 共玩 $600$ 局，预计约 $100$ 次掷出六点，$100$ 次掷出五点，$400$
次掷出其他点数。奖金约为 $100 \times 6 + 100 \times 3 = 900$。成本：
$600 \times 2 = 1200$。期望净收益为 $900 - 1200 = -300$，即
每局平均亏损约 $0.5$ ——这个游戏
不值得玩。
\end{solution}

\begin{solution}{pb:g10:proba:1}
\textbf{1.} $P(\text{国王或红心}) = \frac{4}{52} +
\frac{13}{52} - \frac{1}{52} = \frac{16}{52} = \frac{4}{13}$：
红心国王同时属于这两个事件，不应计算两次，
因此需要做减法
（\cref{thm:g10:proba:addition}）。

\textbf{2.} $P(\text{没有正面}) = \left(\frac12\right)^3 =
\frac18$，故 $P(\text{至少一次正面}) = \frac78$。

\textbf{3.} $P(RR) = \frac35 \times \frac24 = \frac{3}{10}$；
$P(\text{一红一蓝}) = \frac35 \times \frac24 + \frac25
\times \frac34 = \frac{6}{20} + \frac{6}{20} = \frac35$。

\textbf{4.} $P(BB) = \frac25 \times \frac14 = \frac1{10}$；而
$\frac{3}{10} + \frac{6}{10} + \frac{1}{10} = 1$：树状图的
末端概率之和恒为 $1$。

\textbf{5.} 互斥事件满足 $P(A \cup B) = P(A) +
P(B) = 1.1 > 1$：不可能 --- 有某个概率是凭空捏造的。

\textbf{6.} （多数人，包括那一千位写信者，
回答“无所谓，五五开”。请保留你对
第 8 题的答案。）

\textbf{7.} 汽车在 1 号门后（概率 $\frac13$）：主持人
打开任一扇山羊门；换选则会失去汽车：\emph{输}。
汽车在 2 号门后（$\frac13$）：主持人\emph{必须}打开 3 号门；
换选则会选中 2 号门：\emph{赢}。汽车在 3 号门后：
情形对称：\emph{赢}。总计：
$P(\text{换选获胜}) = \frac13 + \frac13 = \frac23$。

\textbf{8.} 你第一次选择出错的概率为
$\frac23$ --- 恰在此时，主持人被迫的揭示
使汽车留在剩下的那扇门后：换选会把
最初的每一次选错都变成胜利。$\frac23$，
无论直觉怎么说。

\textbf{9.} 六种等可能的情形：汽车在 $1$、$2$
或 $3$ 号门后，主持人打开 $2$ 或 $3$ 号门（绝不是你选的门）。揭示
的结果表明汽车在其中两种情形中（汽车在 2 号门后、开 2 号门，汽车在 3 号门后、开 3 号门）：
这两种被舍弃。剩下的四种情形为：汽车在 1 号门后、开 2 号门，
汽车在 1 号门后、开 3 号门，汽车在 2 号门后、开 3 号门，汽车在 3 号门后、开 2 号门 --- 换选获胜的
是后两种：$\frac24 = \frac12$。主持人若不知情，
五五开\emph{确实}正确：悖论完全存在于
主持人的知情状况，它不均匀地筛选了各种情形。

\textbf{10.} 你选中汽车的概率为
$\frac{1}{100}$；知情主持人揭开的 $98$ 扇门把
其余的 $\frac{99}{100}$ 全部汇集到剩下那扇关着的门上。换选，
则获胜 $99$ 次，共 $100$ 局。

\textbf{11.} 三张牌（一张 A、两张王），一位朋友
偷看之后，总从你未选的两张牌中弃掉一张王；
你玩 $300$ 局，每局都换选，并记录获胜次数。
期望频率为 $\frac23 \approx 0.667$，波动范围
为 $\pm\frac{1}{\sqrt{300}} \approx 0.058$，预计换选获胜的比例约在
$61\,\%$ 至 $72\,\%$ 之间 --- 远离任何
“五五开”的说法。

\textbf{12.} 等可能的情形（年长者--年幼者）：
“至少有一个男孩”保留 BB、BG、GB：
$P(\text{BB}) = \frac13$。

\textbf{13.} “年长者是男孩”保留 BB、BG：
$P(\text{BB}) = \frac12$。不同的信息会筛选出
不同的情形集合 --- 正如两个蒙蒂所做的那样：你所了解的内容
很重要，而\emph{如何}得知它也同样
重要。

\textbf{14.} 全部相同：HHH、TTT：$\frac28 = \frac14$。
达朗贝尔的错误：他的两种“情形”并不等可能
--- “不全相同”包含了八种结果中的六种。
\Cref{prop:g10:proba:uniform} 只在等可能的
结果上计数。

\textbf{15.} 每局的平均所得：$6 \times \frac16 + 3
\times \frac16 = 1.50$ 欧元，相对于 $2$ 欧元的投注，
平均每局亏损 $0.50$ 欧元，约 $300$ 欧元（
$600$ 局）。公平保费使平均赔付与之相等：
$1\,\% \times 10\,000 = 100$ 欧元；实际保费还需计入成本
与安全余量 --- 保险公司必须经得起糟糕的年份，
而不只是平均的年份。

\textbf{16.} 区间：$0.5 \pm \frac{1}{\sqrt{100}} =
\intcc{0.4}{0.6}$。$55$ 次正面（$0.55$）：在区间内，并不异常。
$70$ 次正面（$0.70$）：远在区间之外 --- 应当怀疑这枚硬币。

\textbf{17.} $\frac{1}{\sqrt{1000}} \approx 0.032$：一位
真实支持率为 $50\,\%$ 的候选人，民调结果通常会落在
$\intcc{0.468}{0.532}$ 之间 --- 而 $0.52$ 稳稳地位于其中。
民调\emph{暗示}出领先，却什么也证明不了：误差范围
足以淹没这一领先。

\textbf{18.} 第一份样本：区间 $0.04 \pm \frac{1}{20} =
\intcc{0}{0.09}$；观测值 $0.07$ 在区间内：无法下结论。
第二份样本：$0.04 \pm \frac{1}{50} = \intcc{0.02}{0.06}$；
观测值 $0.07$ 落在区间外：该论断现在受到严重
怀疑。频率相同，但仪器更精密：精度随
$\sqrt n$ 增长。

\textbf{19.} 波动按 $\frac{1}{\sqrt n}$ 的规律缩小：
样本扩大为四倍，噪声减半。偏差 --- 即询问了错误的人群 ---
无论 $n$ 多大都无法消除：抽样不当的两百万张选票
结果依然错误（初中分册中的“数据如何说谎”周末问题）。

\textbf{20.} (i) 计数前先检查各结果是否等可能，
达朗贝尔正是忘了这一点。(ii) 树状图的概率沿分支相乘；
“至少一次”则应转而求对立事件。
(iii) 信息会筛选情形，而信息的\emph{来源}
会改变筛选方式 --- 对比知情的蒙蒂与偶然开门的蒙蒂，
对比“年长者是男孩”与“至少有一个男孩”。(iv) 频率以
$\frac{1}{\sqrt n}$ 的速度趋近概率，因此一次耐心的模拟
足以平息任何争论 --- 包括与一万名
愤怒来信者之间的争论。
\end{solution}
